\documentclass[10pt]{article}

\usepackage[margin=0.75in]{geometry}
\usepackage{amsmath,amsthm,amssymb,bm}
\usepackage{xcolor}
\usepackage{cancel}
\usepackage{graphicx}
\usepackage{changepage}
\usepackage{circuitikz}
\usepackage{pgfplots}
\usepackage{minted}
\usepackage{physics}
\usepackage{tikz}
\usepackage{tikz-3dplot}
\usepackage{hyperref}
\usepackage{siunitx}
\usepackage{fontspec}
\usepackage{relsize}
\usepackage{subfig}
\usepackage{todonotes}
\usepackage{contour}
\usepackage{multicol, multirow, booktabs}
\usepackage[breakable]{tcolorbox}
\usepackage[inline]{enumitem}

\theoremstyle{definition}
\newtheorem{problem}{Problem}
\newtheorem{soln}{Solution}

\pgfplotsset{compat=newest}
\usetikzlibrary{lindenmayersystems}
\usetikzlibrary{arrows}
\usetikzlibrary{calc}
\usetikzlibrary{positioning, fit}
\usetikzlibrary{3d, perspective}
\usetikzlibrary{math}
\usetikzlibrary{fadings}
\usetikzlibrary{angles,quotes} 

\definecolor{incolor}{HTML}{303F9F}
\definecolor{outcolor}{HTML}{D84315}
\definecolor{cellborder}{HTML}{CFCFCF}
\definecolor{cellbackground}{HTML}{F7F7F7}
\newcommand{\ui}{\hat{i}}
\newcommand{\uj}{\hat{j}}
\newcommand{\uk}{\hat{k}}
\newcommand{\ux}{\hat{\mathbf{x}}}
\newcommand{\uy}{\hat{\mathbf{y}}}
\newcommand{\uz}{\hat{\mathbf{z}}}
\newcommand{\uv}[1]{\hat{\bv{{#1}}}}
\newcommand{\pr}[1]{{#1^\prime}}
\newcommand{\pd}[2][]{{\frac{\partial^{#1}}{\partial {{#2}^{#1}}}}}
\newcommand{\dif}[2][]{{\frac{d^{#1}}{d{{#2}^{#1}}}}}
\newcommand{\grd}{\nabla}

\pgfdeclarelayer{background}  
\pgfsetlayers{background,main}
\AtBeginDocument{\RenewCommandCopy\qty\SI}

\makeatletter
\newcommand{\boxspacing}{\kern\kvtcb@left@rule\kern\kvtcb@boxsep}
\makeatother
\newcommand{\prompt}[4]{
    \ttfamily\llap{{\color{#2}[#3]:\hspace{3pt}#4}}\vspace{-\baselineskip}
}

\newcommand{\thevenin}[2]{
  \begin{center}
    \begin{circuitikz} \draw
      (0,0) -- (2,0) to[battery1, l_=$V_{Th}\eq#1$] (2,2) 
      to[resistor, l_=$R_{Th}\eq#2$] (0,2)
      ;
      \draw [o-] (-.07,2.079);
      \draw [o-] (-.07,0.079);
    \end{circuitikz}
  \end{center}
}

\newcommand{\norton}[2]{
  \begin{center}
    \begin{circuitikz} \draw
      (0,0) -- (3,0) to[american current source, l_=$I_{N}\eq#1$] (3,2) -- (0,2) (2,0)
      to[resistor, l=$R_{N}\eq#2$] (2,2)
      ;
      \draw [o-] (-.07,2.079);
      \draw [o-] (-.07,0.079);
    \end{circuitikz}
  \end{center}
}

\DeclareMathOperator{\Div}{div}
\DeclareMathOperator{\Curl}{curl}
\DeclareMathOperator{\sgn}{sgn}

\newcommand{\highlight}[1]{\colorbox{yellow}{$\displaystyle #1$}}

\newcommand{\ti}[1]{\widetilde{#1}}

\newfontface{\Kaufmann}{Kaufmann}
\DeclareTextFontCommand{\kf}{\Kaufmann}
\newcommand{\scriptr}{\fontsize{12pt}{12pt}\kf{r}\fontsize{10pt}{10pt}}

\newfontface{\KaufmannB}{Kaufmann Bold}
\DeclareTextFontCommand{\kfb}{\KaufmannB}
\newcommand{\bscriptr}{\fontsize{12pt}{12pt}\kfb{r}\fontsize{10pt}{10pt}}
\newcommand{\justif}[2]{&{#1}&\text{#2}}
\newcommand{\bv}[1]{\bm{\mathrm{#1}}}

\title{Physics 3200Y: Assignment VIII}
\author{Jeremy Favro (0805980) \\ Trent University, Peterborough, ON, Canada}
\date{\today}

\begin{document}
\maketitle

% PROBLEM 1
\begin{problem}
A rotating disc (radius $R$, angular velocity $\omega$) has a uniform surface charge density $\sigma$. Find:
\begin{enumerate}[label=(\alph*)]
  \item The surface current density $\bv{K}$.
  \item The magnetic field at a point $z > R$ along the axis of symmetry. Note the similarity to
        Example 5.6.
  \item The magnetic dipole moment of the disc. You can do this by treating the disc as a sum
        of circular wire loops, and adding up the magnetic dipole moments.
  \item Now compare the magnetic field due to the dipole you found in (c) to the magnetic field
        you found in (b). They should be the same in the limit $z \gg R$. To do the comparison,
        you'll need to evaluate the Taylor series for $(1 + X)^{-1/2}$ to second order in $X$. Do this
        explicitly and find the leading-order nonzero term.
\end{enumerate}
\end{problem}
\begin{soln}~
  \begin{enumerate}[label=(\alph*)]
    \item A surface current density is by definition the current per unit width. If we envision the rotating disc as a bunch individual
          wire loops each carrying a piece of the charge $\sigma$ at radius $s'$ and tangential velocity $\bv{\omega}\times \bv{s'}$ we see that the current is
          $\bv{K}=\sigma \omega s'\uv{\phi}=\sigma \omega s'\left(-\sin\phi\ux+\cos\phi\uy\right)$
    \item Here $\bv{r}=(0,0,z)$ and $\bv{r}'=(x',y',0)$ so $\bscriptr=(-x',-y',z)$ and when parametrizing the disc in $s'$ and $\phi'$, $\bscriptr=(-s'\cos\phi',-s'\sin\phi', z)$.
          We want to use
          $$\bv{B}(\bv{r})=\frac{\mu_0}{4\pi}\int da' \frac{\bv{K}(\bv{r}')\times \uv{\bscriptr}}{\scriptr^2}$$
          so we compute
          \begin{align*}
            \bv{K}\times \bscriptr =\begin{vmatrix}
                                      \ux                       & \uy                      & \uz \\
                                      -\sigma \omega s'\sin\phi & \sigma \omega s'\cos\phi & 0   \\
                                      -s'\cos\phi'              & -s'\sin\phi'             & z
                                    \end{vmatrix}
             & =\sigma z\omega s'\cos\phi\ux+\sigma z \omega s'\sin\phi\uy + \left[\sigma \omega s'^2\sin^2\phi+\sigma \omega s'^2\cos^2\phi\right]\uz \\
             & =\sigma z\omega s'\cos\phi\ux+\sigma z \omega s'\sin\phi\uy + \sigma \omega s'^2\uz
          \end{align*}
          so the integral becomes
          \begin{align*}
            \bv{B}(\bv{r}) & =\frac{\mu_0}{4\pi}\int da' \frac{\sigma z\omega s'\cos\phi\ux+\sigma z \omega s'\sin\phi\uy + \sigma \omega s'^2\uz}{\scriptr^3}                                     \\
                           & =\frac{\mu_0}{4\pi}\int_{0}^{2\pi}\int_{0}^{R} \frac{\sigma z\omega s'\cos\phi\ux+\sigma z \omega s'\sin\phi\uy + \sigma \omega s'^2\uz}{(s'^2+z^2)^{3/2}}s'ds'd\phi' \\
                           & =\frac{\mu_0 \sigma \omega }{2}\int_{0}^{R} \frac{s'^3}{(s'^2+z^2)^{3/2}}ds'\uz
          \end{align*}
          here we let $u=s'^2+z^2\implies du =2s'ds'$ so that
          \begin{align*}
            \bv{B}(\bv{r}) & =\frac{\mu_0 \sigma \omega }{2}\int_{0}^{R} \frac{s'^3}{(u)^{3/2}}\frac{du}{2s'}\uz                                    \\
                           & =\frac{\mu_0 \sigma \omega }{4}\int_{0}^{R} \frac{u-z^2}{(u)^{3/2}}du\uz                                               \\
                           & =\frac{\mu_0 \sigma \omega }{4}\left[\int_{0}^{R} \frac{u}{(u)^{3/2}}du-\int_{0}^{R} \frac{z^2}{(u)^{3/2}}du\right]\uz \\
                           & =\frac{\mu_0 \sigma \omega }{4}\left[2\sqrt{u}+2\frac{z^2}{\sqrt{u}}\right]_{0}^{R}\uz                                 \\
                           & =\frac{\mu_0 \sigma \omega }{2}\left[\frac{R^2+2z^2}{\sqrt{R^2+z^2}}-\frac{2z^2}{\sqrt{z^2}}\right]\uz                 \\
                           & =\frac{\mu_0 \sigma \omega }{2}\left[\frac{R^2+2z^2}{\sqrt{R^2+z^2}}-2z\right]\uz
          \end{align*}
    \item Doing this as the question suggests: we know that the magnetic moment of a wire loop of radius $s'$ is $d\bv{m}=dI\pi s'^2\uz$.
    We don't know $dI$ so we'll have to do some work to find it. We do know $\bv{K}=\sigma \omega s'\uv{\phi}$ for the entire surface 
    and we know that $$\bv{K}=\frac{d\bv{I}}{d\ell}$$
    so $$\int \bv{K} d\ell = \bv{I}=\frac{\sigma \omega s'^2}{2}\uv{\phi}\implies $$
          % Here we only know $\bv{K}$ so we'll have to integrate to find the current of each loop. Since $\bv{K}=\frac{d\bv{I}}{d\ell}$,
          % $$I=\int_\mathcal{C} \bv{K}\cdot d\bv{\ell}$$
          % and here $\bv{K}\cdot d\bv{\ell}=\sigma \omega ds'\uv{\phi}\cdot s'd\phi\uv{\phi}=\sigma \omega s'd\phi ds'$.
          % Note the $ds'$ in $\bv{K}$. This is because we are not considering $\bv{K}$ for the whole surface but instead a slice of it given by the wire loop of
          % thickness $ds'$. This gives
          % $$I=\int_0^{2\pi}\sigma \omega s' d\phi ds'=2\pi\sigma\omega s'ds'$$
          % This means that
          % $$d\bv{m}=2\pi^2\sigma\omega s'^3ds' \uz.$$
          % Now since we are treating the disc as a continuous collection of thin wire loops of radius $0\to R$ we integrate the above expression
          % along $s'$ to find the total moment,
          % $$\bv{m}=\int_{0}^{R}2\pi^2\sigma \omega s'^3ds'\uz=\frac{\pi^2\sigma \omega R^4}{2}\uz.$$
    \item We know the magnetic field of a dipole to be
          $$\bv{B}_{dip}(\bv{r})=\frac{\mu_0}{4\pi}\left[\frac{3(\bv{m}\cdot \bv{r})\bv{r}}{r^5}-\frac{\bv{m}}{r^3}\right].$$
          Here $\bv{r}=(0,0,z)$ so
          $$\bv{B}_{dip}(\bv{r})=\frac{\mu_0}{4\pi}\left[\frac{3(\frac{\pi^2\sigma \omega R^4}{2})z\uz}{z^{5/2}}-\frac{\frac{\pi^2\sigma \omega R^4}{2}\uz}{z^{3/2}}\right]
            =\frac{\mu_0}{4\pi}\left[\frac{\pi^2\sigma \omega R^4}{z^{3/2}}\right]\uz
            .$$
          Now the ``true'' magnetic field is given by
          $$\bv{B}(\bv{r})=\frac{\mu_0 \sigma \omega }{2}\left[\frac{R^2+2z^2}{\sqrt{R^2+z^2}}-2z\right]\uz$$
          which for $z\gg R$ becomes
          \begin{align*}
            \bv{B}(\bv{r}) & =\frac{\mu_0 \sigma \omega }{2}\left[\left[R^2+2z^2\right]\left[R^2+z^2\right]^{-1/2}-2z\right]\uz \\
                           & =\frac{\mu_0 \sigma \omega }{2}\left[R\left[R^2+2z^2\right]\left[1+\left(\frac{z}{R}\right)^2\right]^{-1/2}-2z\right]\uz \\
                           & =\frac{\mu_0 \sigma \omega }{2}\left[R\left[R^2+2z^2\right]\left[1-\frac{1}{2}\left(\frac{z}{R}\right)^2\right]-2z\right]\uz \\
                           & =\frac{\mu_0 \sigma \omega }{2}\left[\left[2z^2\right]\left[R-\frac{1}{2}\frac{z^2}{R}\right]-2z\right]\uz \\
                           & =\frac{\mu_0 \sigma \omega }{2}\left[2Rz^2-\frac{z^4}{R}-2z\right]\uz \\
          \end{align*}
  \end{enumerate}
\end{soln}

% PROBLEM 2
\begin{problem}
In order to obtain Eq. (6.15) of Griffiths, we needed to show that
$$\int_\mathcal{V}\bv{\grd}' \times \left[\frac{\bv{M}(\bv{r}')}{\scriptr}\right]d^3r'=-\oint_\mathcal{S}\frac{\bv{M}(\bv{r}')\times d\bv{a}'}{\scriptr}$$
Complete problem 1.61(b), and then use this to fill in the missing steps of the derivation of
Eq. (6.12).
\end{problem}
\begin{soln}
Problem 1.61(b) asks for a proof of 
$$\int_{\mathcal{V}}\bv{\grd}\times \bv{v}d\tau = -\oint_\mathcal{S}\bv{v}\times d\bv{a}.$$
The problem suggests replacing $\bv{v}$ by $\bv{v}\times \bv{c}$ in the divergence theorem, so that is what I will do.
The divergence theorem says that 
$$\int_{\mathcal{V}}\bv{\grd}\cdot \bv{v} d\tau=\oint_\mathcal{S}\bv{v}\cdot \uv{n} da$$
and doing the suggested replacement gives us 
$$\int_{\mathcal{V}}\bv{\grd}\cdot (\bv{v}\times \bv{c}) d\tau=\oint_\mathcal{S}(\bv{v}\times \bv{c})\cdot \uv{n} da.$$
Using product rule (6) from the front cover of Griffiths and the fact that $\uv{n} da =d\bv{a}$ gives us 
$$\int_{\mathcal{V}} \bv{c}\cdot (\bv{\grd}\times \bv{v})-\bv{v}\cdot (\bv{\grd}\times \bv{c}) d\tau=\oint_\mathcal{S}(\bv{v}\times \bv{c})\cdot d\bv{a}.$$
Now we use triple product rules (1) from the front cover of Griffiths to get that 
$$\int_{\mathcal{V}} \bv{c}\cdot (\bv{\grd}\times \bv{v})-\bv{v}\cdot (\bv{\grd}\times \bv{c}) d\tau=-\oint_\mathcal{S}\bv{c}\cdot (\bv{v}\times d\bv{a})$$
where we've flipped the sign as $\bv{a}\times \bv{b} = -\bv{b}\times \bv{a}$. Now note that (and I assume this is why $\bv{c}$ is called what it is) for constant
$\bv{c}$, $\bv{\grd}\times \bv{c}=0$ so 
$$\bv{c}\cdot \int_{\mathcal{V}} (\bv{\grd}\times \bv{v})d\tau=-\bv{c}\cdot \oint_\mathcal{S}(\bv{v}\times d\bv{a})$$
which only implies that the components of these integrals along $\bv{c}$ are equal, not that the integrals themselves are equal. I.e.
consider $\bv{a}=\ux$, $\bv{b}=x\ux+y\uy$, $\bv{c}=x\ux+8y\uy+123zyx\uz$
then $\bv{a}\cdot\bv{b}=\bv{a}\cdot\bv{c}$ certainly but $\bv{b}\neq\bv{c}$.
\end{soln}

% PROBLEM 3
\begin{problem}
Griffiths problem 6.1.
\end{problem}
\begin{soln}

\end{soln}

% PROBLEM 4
\begin{problem}
Griffiths problem 6.7.
\end{problem}
\begin{soln}

\end{soln}
\end{document}